\subsection{内射模的结构}

引理4.——设M是非零诺特A-模，并设I是M的内射包络。则I具有一个不可分解的内射子模。

我们显然可以假设M是I的子模。设N是M的子模，使得I不是N的内射包络，并且N对此性质是极大的。由注记2（X，第21页），存在I的子模I₁，它是N的内射包络；则I₁是I的一个直和项，并设J为其补。我们有J ≠ 0；让我们证明J是不可分解的。若J'是J的非零直和项子模，则J' ∩ M ≠ 0，且

\[
\left(\mathrm{J} ^ {\prime} \cap \mathrm{M}\right) \cap \mathrm{N} \subset \mathrm{J} ^ {\prime} \cap \mathrm{I} _ {1} = 0.
\]

子模 \(N' = (J' \cap M) + N\) 是M的子模，它是 \(J' \cap M\) 与N的直和，因此严格包含N。此外， \(N'\) 包含于J的子模 \(J' + I_{1}\) 中，而该子模是 \(J'\) 和 \(I_{1}\) 的直和，因而是内射的。由N的极大性，这意味着 \(J' + I_{1} = I\) ，因此 \(J' = J\) ，且J是不可分解的。

我们用 \(\mathcal{I}\) 表示不可分解内射A-模的类所成的集合（X，第21页，推论1）。

回顾（X，第17页，推论1），若A是左诺特的，则由内射A-模组成的每个直和都是内射A-模。

定理3. --- 设I是一个内射A-模。

\begin{enumerate}
\def\labelenumi{\alph{enumi})}
\item
  若 I 是诺特 A-模 M 的内射包络，则 I 是有限个不可分解（内射）子模的直和。
\item
  若A是左诺特的，I是一族不可分解（内射）子模的直和。
\item
  如果 I 是（内射）不可分解子模的直和，则存在唯一的一族基数 \((a_{\mathbf{E}})_{\mathbf{E}\in\mathcal{I}}\) ，使得 I 同构于
\end{enumerate}

\[
\bigoplus_ {\mathbf {E} \in \mathcal {I}} \mathrm{E} ^ {(a _ {\mathbf {E}})}.
\]

我们首先注意到，c) 可由命题14（X，第21页）以及第八章第1节第7号定理2推出。下面证明a)。

设 N 为 M 的一个子模，其内射包络是有限个不可分子模的直和，且 N 对此性质是极大的（这样的 N 存在，因为 M 是诺特的）。由注记 2（X，第 21 页），存在 I 的一个子模 I \(_{1}\) ，它是N的内射包络。如果 I \(_{1}\) = I，则证明完毕；否则，令J为I₁在I中的补子模。则J是Noether模J ∩ M的内射包络（同上所引），因此具有一个不可分解的内射子模J'（引理4）。于是I₁ + J'是内射的，是有限个不可分解子模的直和，并且是M的子模(I₁ + J') ∩ M的内射包络，该子模严格大于N，从而得出矛盾。

假设A是左诺特环，我们来证明 \(b\) )。设 X 为诸集合 \(\mathrm{Hom}_{\mathbf{A}}(\mathrm{E},\mathrm{I})\) （ \(\mathrm{E}\in\mathcal{I}\) ）的不交并。对 X 的每个子集 Y，赋予一个 A-模 \(\mathrm{E}_{\mathbf{Y}}\) 以及一个 A-同态 \(f_{\mathbf{Y}}:\mathrm{E}_{\mathbf{Y}}\to\mathrm{I}\) 如下：Y 是某个族 \((\mathrm{Y}(\mathrm{E}))_{\mathrm{E}\in\mathcal{I}}\) 其中 \(\mathrm{Y}(\mathrm{E})\subset\mathrm{Hom}_{\mathbf{A}}(\mathrm{E},\mathrm{I})\) ，我们设

\[
\mathrm{E} _ {\mathbf {Y}} = \bigoplus_ {\mathbf {E} \in \mathcal {I}} \mathrm{E} ^ {(\mathbf {Y} (\mathbf {E}))}
\]

以及 \(f\) 在直和项 \(\mathbf{E}_{\mathbf{Y}}\) 上对应于元素 \(y\) （属于 \(\mathbf{Y}(\mathbf{E}) \subset \operatorname{Hom}_{\mathbf{A}}(\mathbf{E}, \mathbf{I})\) ）的分量是 \(y: \mathbf{E} \to \mathbf{I}\) 。设 Y 为 X 的一个子集，使得 \(f_{\mathbf{Y}}\) 是单射且 Y 关于此性质是极大的（这样的子集由 E, III, p.～20 存在）；只需证明 \(f_{\mathbf{Y}}\) 是双射。否则，设 J 为内射子模 \(\operatorname{Im}(f_{\mathbf{Y}})\) 在 I 中的一个补子模；由于 J 非零，它有一个非零的 Noether 子模（因为 A 被假定为 Noether 的），从而也有一个非零的内射子模 J'，它是某个 Noether 模的内射包络。由 a)，J' 是有限个非零不可分子模的直和。因此存在 X 的一个非空有限子集 Y，使得 \(f_{\mathbf{Y}}\) 将 \(\mathbf{E}_{\mathbf{Y}'}\) 双射地映到 J' 上。由于 \(\operatorname{Im}(f_{\mathbf{Y}}) \cap \mathbf{J}' = 0\) ，我们有 \(\mathbf{Y} \cap \mathbf{Y}' = \varnothing\) 且 \(f_{\mathbf{Y} \cup \mathbf{Y}'}\) 是单射；这与 Y 的极大性矛盾，从而完成了证明。

